\section{From the codebook to a served file}
\label{sec:model}

A served file holds more than lattice codes. This section lists what else it
holds and describes the three changes that make up the files we measure. None
of them changes a lattice code or the decoder.

\subsection{What a file holds, and how we count it}
\label{sec:file}

Qwen3~\citep{qwen3_2025} has seven projection matrices per layer ($q$, $k$,
$v$, $o$, gate, up, down): 252 in the 36 layers of the 4B and 8B models, 280 in
the 40 layers of the 14B. Each projection is stored either as \Tetra{} blocks
or as an int4 record: 4-bit weights in groups of 128 with an f16 scale and
offset per group, read by a separate kernel. A \Tetra{} projection also keeps a
small tail of weights exactly, in f16, and one scale per output row. The norms
stay in f16. The embedding table is stored apart, and so is the output head
when it is not tied to the embedding, which is the case at 8B and 14B.

We count memory in bits per parameter over the whole model: every byte the
served model keeps for its weights, including the embedding, the tail, the
scales and the int4 records, divided by the number of parameters. One tool
computes it from the file for every object in this paper. The KV cache and the
activations are not counted.

\subsection{Training the row scales}
\label{sec:rowscales}

Equation~\eqref{eq:recon} ends with one scale per output row,
$\sigma_{\mathrm{row}}$, fitted by least squares during encoding. We retrain
these scales and nothing else, with a KL loss against the FP16 model as
teacher~\citep{hinton2015} on DCLM-edu text~\citep{dclmedu}. The lattice codes
stay fixed and no bit is added: the trained scales are folded back into the
same field. At 4B the run sees 19.5 million tokens and takes 1.7~h on one
L40S; the 8B and 14B runs take 1.1~h and 2.0~h on one H200. This step alone
adds 3.15, 3.29 and 1.67 MMLU points at 4B, 8B and 14B (\S\ref{sec:chain}).

\subsection{int4 where the lattice loses the most}
\label{sec:int4}

The projections do not all lose the same. Every file stores \texttt{v\_proj}
in int4: under grouped-query attention these matrices are small, and at 4B
moving them from \Tetra{} to int4 gained 3.47 MMLU points on the 2{,}280-question
sample for 0.05 bits per parameter. The files we call \emph{sealed} go further
and also store \texttt{o\_proj} and a window of middle layers of
\texttt{down\_proj} in int4 (Table~\ref{tab:files}). At 8B only
\texttt{down\_proj} is added; \S\ref{sec:chain} shows why.

\subsection{4-bit embedding tables}
\label{sec:q4emb}

Qwen3's vocabulary has 151{,}936 entries, so the embedding table is 9.7\,\% of
the parameters at 4B. At 8B and 14B the output head is a second table of the
same shape, and the two together are 15.2\,\% and 10.5\,\% of the parameters.
Our earlier files stored these tables in 8 bits. The sealed files store them in
4 bits, in groups of 64 with an f16 scale and offset per group, the scheme MLX
uses~\citep{mlx2023}. The GPU dequantizes rows on the fly, in the gather at the
input and inside the matrix-vector product of the head. Going from 8 to 4 bits
frees 0.19~GB at 4B, 0.62~GB at 8B and 0.78~GB at 14B. That is what pays for
the int4 projections: every sealed file needs fewer bits per parameter than the
file it was built from (\S\ref{sec:chain}).

\begin{table}[t]
\centering\small
\begin{tabular}{lrrr}
\toprule
 & Qwen3-4B & Qwen3-8B & Qwen3-14B\\
\midrule
projections, \Tetra{} $+$ int4 & 168 $+$ 84 & 199 $+$ 53 & 181 $+$ 99\\
int4 projections & $v$, $o$, down 12--23 & $v$, down 10--26 & $v$, $o$, down 10--28\\
4-bit embedding tables & 1 (tied head) & 2 & 2\\
file size (bytes) & 1{,}418{,}224{,}685 & 2{,}815{,}098{,}745 & 5{,}087{,}000{,}541\\
bits per parameter, whole model & 2.7320 & 2.6953 & 2.7305\\
\bottomrule
\end{tabular}
\caption{The three sealed files. Layers are numbered from 0. Bits per
parameter count every weight byte of the served model, embedding included, as
defined in \S\ref{sec:file}. Data: paper2-chain.csv.}
\label{tab:files}
\end{table}
